\documentclass[11pt]{amsart}
\usepackage[T1]{fontenc}
\usepackage{lmodern}
\usepackage[margin=1in]{geometry}
\usepackage{amsmath,amssymb,amsthm,booktabs,array,xcolor}
\usepackage{microtype}
\usepackage[colorlinks=true,linkcolor=blue!45!black,citecolor=blue!45!black,urlcolor=blue!45!black]{hyperref}
\newtheorem{theorem}{Theorem}[section]
\newtheorem{lemma}[theorem]{Lemma}
\newcommand{\Z}{\mathbb Z}
\newcommand{\artifact}[2]{\href{https://github.com/shjvandekamp/math/blob/main/topics/snaky/#1}{#2}}
\newcommand{\snakycell}[2]{%
  \put(#1,#2){\color{blue!12}\rule{\unitlength}{\unitlength}}%
  \put(#1,#2){\framebox(1,1){}}}
\hypersetup{pdftitle={Snaky within nineteen Maker moves},
  pdfsubject={A finite certificate and Lean kernel proof for the Maker--Breaker game on the infinite square board}}
\title[Snaky within nineteen Maker moves]{Snaky within nineteen Maker moves}
\date{15 October 2026}
\begin{document}
\begin{abstract}
We give a finite, independently checked certificate that Maker can force
Snaky within nineteen of Maker's own moves on the initially empty infinite
square board. The same certificate proves the bound on a $17\times17$ board.
It combines N\'andor Sieben's twenty-move strategy with local tactics from
Trevor Morris's $15\times15$ certificate. The construction retains all
acquired Maker stones and tracks a larger region known to contain no Breaker
stone, allowing previously discarded safe cells to be used again. The
resulting 3,468-card certificate uses the existing declarative rules. Its
soundness follows by induction on card height, with a universal response
class for moves outside the support. The complete certificate and the
infinite-board soundness theorem have also been checked by the Lean kernel.
We make no optimality claim.
\end{abstract}
\maketitle

\section{Game and theorem}
The board is either $\Z^2$ or $Q=\{0,\ldots,16\}^2$. Maker moves first.
On each turn the player claims exactly one previously unclaimed cell.
Breaker only blocks; Breaker completing a shape does not end the game.
Maker wins upon owning a translated, rotated, or reflected copy of
\begin{equation}\label{eq:snaky}
 F=\{(0,0),(1,0),(2,0),(3,0),(3,1),(4,1)\}.
\end{equation}
\begin{center}
\setlength{\unitlength}{5.5mm}
\begin{picture}(5,2)
\snakycell{0}{0}\snakycell{1}{0}\snakycell{2}{0}\snakycell{3}{0}
\snakycell{3}{1}\snakycell{4}{1}
\end{picture}
\end{center}
Other Maker stones are allowed: the six-cell condition specifies a winning
subset. A move bound counts actual Maker claims, including the winning claim,
and does not count alternating plies.

\begin{theorem}\label{thm:main}
From the empty board, Maker has a strategy that completes a copy of $F$ by
Maker's nineteenth move against every legal Breaker strategy. This holds
both on $\Z^2$ and on $Q$.
\end{theorem}

The data establishing the theorem are
\artifact{certificate/snaky19.cards.txt}{\nolinkurl{certificate/snaky19.cards.txt}},
with SHA-256
\begin{center}\footnotesize\ttfamily
baf8774e1fc6eea60b1756437e8b3caae0ede5e001b10c0c3787aab51840782e
\end{center}
The file contains 3,468 cards, all reachable from root 3467. The root has
height 19, no required Maker stones, first move $(8,8)$, and a 249-cell support
contained in $Q$. Its longest dependency route also has height 19.

\section{Declarative certificate and soundness}
A card is a tuple $C=(A,S,p,h)$: finite integer-cell sets $A$ and $S$, a pivot
$p\in S$, and a positive integer height $h$. Its intended meaning is that
Maker can win within $h$ further claims whenever Maker owns $A$ and Breaker
owns no cell of $S$. Extra Maker stones are permitted. Cards can be placed
by any map $g(x)=Dx+t$, where $D$ is one of the eight symmetries of the square
and $t\in\Z^2$.

A card is terminal if $A\cup\{p\}$ contains a copy of $F$. Otherwise it must
provide the following hints. Each hint selects another card
$(A',S',p',h')$ and a map $g$, with
\begin{equation}\label{eq:hint}
 h'<h,\qquad g(A')\subseteq A\cup\{p\}.
\end{equation}
For every $w\in S\setminus(A\cup\{p\})$, an inside-response hint must
additionally satisfy $g(S')\subseteq S\setminus\{w\}$. An outside-response
hint must satisfy $g(S')\subseteq S$. The latter is one universally valid
rule, not an assumption that Breaker can pass. It applies to every legal
claim outside $S$, including holes in its bounding rectangle and points
arbitrarily far away.

The finite checker verifies these inclusions, all required response classes,
the coordinate maps, terminal geometry, and strict height decreases. The
construction algorithm is not part of the certificate's logical premises.

\begin{lemma}[Soundness]\label{lem:soundness}
Let every card satisfy the preceding rules. At a Maker turn on $\Z^2$, let
$M,W$ be the finite, disjoint ownership sets. If an instance of a height-$h$
card satisfies $g(A)\subseteq M$ and $g(S)\cap W=\varnothing$, Maker can
complete Snaky within at most $h$ additional claims.
\end{lemma}
\begin{proof}
Induct strongly on $h$. A position already containing Snaky is already won.
Otherwise, if $g(p)$ is unowned, Maker claims it; the claim is legal because
$p\in S$ and the support is Breaker-free. If $g(p)$ is already Maker's, Maker
instead claims any fresh cell. Such a cell exists because the histories are
finite and the board is infinite. This is an actual claim and is charged to
the height budget. In either case the resulting Maker set contains
$g(A\cup\{p\})$.

If the terminal condition holds, Maker now owns the required six-cell copy.
Otherwise consider an arbitrary legal Breaker claim $w$. If $w\in g(S)$,
its unique preimage is outside $A\cup\{p\}$, since Breaker cannot claim
Maker's stones. The corresponding checked hint gives a child whose required
stones are owned and whose support avoids both the old Breaker stones and
$w$. If $w\notin g(S)$, the outside hint supplies the same conclusion.
Composition of square symmetries and translations preserves the allowed maps
and winning shapes. The child has height at most $h-1$; induction therefore
bounds the total by one Maker claim plus $h-1$. These two reply cases exhaust
$\Z^2$.
\end{proof}

For the finite board, the same proof uses the capacity invariant
\begin{equation}\label{eq:capacity}
 |M|+|W|+2h\leq |Q|.
\end{equation}
It supplies a free cell for an already-owned pivot's replacement claim and
for a nonterminal Breaker reply. After one claim by each player and a strict
height decrease, the invariant persists. Root support containment keeps all
subsequent supports on the board; child requirements are contained in already
owned stones. At the empty root, $2\cdot19\leq289$. Applying
Lemma~\ref{lem:soundness} directly on the plane, and its finite version on
$Q$, proves both assertions of Theorem~\ref{thm:main}. The infinite-board
conclusion does not rely on treating outside-board Breaker moves as legal
passes in a finite game.

\section{Construction by retaining history and safe space}
The starting twenty-move certificate has 1,738 cards. Reconstructing its
dependency graph gives 15,696 maximum-height labelled routes through 59
critical cards. Full-history replay distinguishes 9,792 routes that repeat
a planned Maker stone, 2,592 with an already-occupied Breaker reply, 864 with
an earlier global win, and 2,448 genuine twenty-move routes. Thus simply
correcting this bookkeeping does not establish nineteen moves.

Fix the initial baseline support $U$. During construction retain the full
Maker set $M$ and the actual Breaker set within $U$, giving the safe region
$E=U\setminus W$. A local template may forget stones and reduce its support,
but these known facts about the actual position remain available.

For a remaining budget $b$, the constructor tries a transformed catalog card
with height at most $b$, requirements contained in $M$, and support contained
in $E$. If one fits, its certified strategy closes that node. Otherwise the
constructor follows the baseline move. Replies whose existing hints already
fit the smaller budget retain those hints; the other replies are processed
recursively. Every claim inside $U$ remains a separate possible update of
$W$, even when outside the current template. Only claims outside $U$ are
merged, since every future support remains inside $U$.

Using only the twenty-move catalog left an unresolved branch. Adding the
8,671-card catalog underlying Morris's $15\times15$ result closed the whole
root. That catalog's root has height 25; its usefulness here lies in its local
tactics, not in a better standalone move bound.

The successful construction uses 180 catalog substitutions. All 180 reclaim
safe cells outside the current baseline support; 68 also use Maker stones
forgotten by the baseline requirements. Seventeen distinct $15\times15$
tactics account for 164 substitutions, and three twenty-move tactics account
for the remaining 16. No zero-cost repeated-pivot shortcut is used in this run.

Three early substitutions illustrate how the surviving long branches are
removed. IDs below are zero-based IDs in the public $15\times15$ catalog.
Coordinates are in the final root frame, immediately before Maker's next claim.
\begin{enumerate}
\item With $M=\{(8,8),(8,9)\}$ and $W=\{(7,8),(9,9)\}$, replace baseline
card 1702 of height 18 by source card 28 of height 17 under
\[
 g(x,y)=(y+2,x+3).
\]
Its next move is $(8,10)$. Two Maker claims have already occurred, so this
continuation has total bound 19.
\item With $M=\{(8,8),(8,9),(11,8)\}$ and
$W=\{(7,8),(8,10),(12,8)\}$, replace baseline card 1036 of height 17 by
source card 79 of height 15 under
\[
 g(x,y)=(14-x,y+3).
\]
Its next move is $(11,9)$, giving total bound 18 on this branch.
\item For the same $M$ and $W=\{(7,8),(8,10),(13,8)\}$, use source card
145 of height 14 under
\[
 g(x,y)=(y+4,12-x).
\]
Its next move is $(9,8)$, giving total bound 17 on this branch.
\end{enumerate}

These examples do not replace verification of the remaining replies. The
constructor exports every new decision as an ordinary card, initially with
requirements $M$ and support $E$. It then discards support cells unused by
its pivot, requirements, and child supports. Reachability pruning and
reindexing leave 1,662 twenty-move-source cards, 1,653 $15\times15$-source
cards, and 153 newly composed cards. Exact source IDs and transformation
witnesses are recorded in the
\artifact{constructor/README.md}{constructor documentation}.

\section{Computational evidence and reproducibility}
The independent Python verifier shares no constructor checking helpers. For
the exact delivered file it checked 61,040 inside-response obligations,
3,456 universal outside-response obligations, and twelve terminal cards.
It accepted empty root requirements, support containment, and height 19.

The verifier's tests include malformed terminal shapes, missing inside or
outside replies, invalid transforms, broken ownership/support inclusions,
invalid references, and incorrect height descent. Controls applied to the
final file rejected both lowering the root height to 18 and deleting its
reply at $(0,4)$. The first rejection concerns this certificate's declared
dependencies; it does not prove that an eighteen-move strategy is impossible.

The deterministic construction used Python 3.9.6, one computational thread,
no random seed, a 120-second wall limit, and a two-million-call cap. The
successful run took about fifteen seconds. The certificate can be checked
offline using only the Python standard library:
\begin{verbatim}
python3 topics/snaky/reproduce.py
\end{verbatim}
Adding \texttt{--reconstruct} regenerates the certificate from the pinned
source files and requires byte-for-byte identity. The final reproduction
passed all 31 tests and both corruption controls, and regenerated identical
bytes.

Lean 4.34.1 completed the full kernel check of the generated certificate and
the direct plane soundness development in
\artifact{formal/Snaky/Plane.lean}{\nolinkurl{formal/Snaky/Plane.lean}}.
The file \artifact{formal/Snaky/Result19.lean}{\nolinkurl{formal/Snaky/Result19.lean}}
proves
\begin{quote}\small\ttfamily
Snaky.snaky\_wins\_17x17\_in\_19\par
Snaky.snaky\_wins\_plane\_in\_19 : Snaky.PlaneBlackWinsIn 19 [] []
\end{quote}
The latter theorem quantifies over every legal integer-coordinate Breaker
reply directly. Both theorems report exactly \texttt{propext},
\texttt{Classical.choice}, and \texttt{Quot.sound} as their axioms. Data checks
use \texttt{decide +kernel}, with no proof placeholders, custom axioms, or
\texttt{native\_decide}. This formal proof does not rely on the correctness of
the Python constructor, verifier, or generator: their output obligations are
checked by the Lean kernel. Exact theorem output and build details appear in
\artifact{formal/kernel-audit.json}{\nolinkurl{formal/kernel-audit.json}};
the formal definitions, Lean kernel, and compiler/runtime remain part of the
formal trust boundary.

\section{Attribution, novelty, and scope}
The twenty-move strategy is due to N\'andor Sieben~\cite{Sieben}, building on
OpenAI's twenty-one-move strategy~\cite{OpenAI}. Trevor Morris's
repository~\cite{Morris} supplies the compact card interface, its soundness
development, the converted twenty-move certificate, and the $15\times15$
tactic catalog. We reuse and credit those tactics; the contribution here is
their history-aware composition into a uniform nineteen-move certificate.
Pinned revisions, hashes, and the retained upstream license are documented in
\artifact{reference/README.md}{\nolinkurl{reference/README.md}}.

The \artifact{NOVELTY.md}{15 October 2026 novelty audit} inspected these
primary sources, the public bounds tracker~\cite{Tracker}, and searches for
nineteen-move-or-better strategies. It located twenty moves as the strongest
public bound and found no matching nineteen-move result. This is a bounded
audit, not proof that no unindexed or unpublished result exists. We establish
an upper bound and claim neither optimality nor any global Breaker lower bound.

\begin{thebibliography}{9}
\raggedright
\bibitem{Sieben}
N\'andor Sieben, \emph{Snaky}, twenty-move proof sequence,
\href{https://github.com/nandorsieben/Snaky/tree/a0a36dd4edffdccdb7c4b14849c81e83dd6a25aa}{\nolinkurl{github.com/nandorsieben/Snaky}},
revision \texttt{a0a36dd4edffdccdb7c4b14849c81e83dd6a25aa}, 8 October 2026.

\bibitem{OpenAI}
OpenAI, \emph{Snaky in 21 Maker moves}, manuscript dated 25 September 2026,
\href{https://github.com/openai/math/tree/1581d4cbdd69a0faf34b2089614fe2a814316a6a}{\nolinkurl{github.com/openai/math}},
repository revision \texttt{1581d4cbdd69a0faf34b2089614fe2a814316a6a},
12 October 2026.

\bibitem{Morris}
Trevor Morris, \emph{Snaky on a $15\times15$ board}, card certificates and Lean
soundness developments,
\href{https://github.com/gotrevor/snaky-15x15/tree/b18ed32473a6a22ebfdda98da62d61b8b1590c1e}{\nolinkurl{github.com/gotrevor/snaky-15x15}},
revision \texttt{b18ed32473a6a22ebfdda98da62d61b8b1590c1e}, 12 October 2026.

\bibitem{Tracker}
Bigospeedrun, public bounds tracker, \nolinkurl{src/lib/other-bounds.json},
\href{https://github.com/bigospeedrun/site/blob/88d1855ba6843586579d8e031389ad94a475bb41/src/lib/other-bounds.json}{\nolinkurl{github.com/bigospeedrun/site}},
revision \texttt{88d1855ba6843586579d8e031389ad94a475bb41}, inspected
15 October 2026.
\end{thebibliography}
\end{document}
